$k_g$を変えても較正順位の仕組みは同じで，変わるのはグラフの密さだけである．NINCO（a）では，$k_g{=}40$の$\sigma$は$r{=}1$から静的$p$を下回る（\NKgNincoKgFourZeroROneMainFpr\%対\NKgNincoKgOneZeroROneStaticFpr\%）が$r{=}30$でも\NKgNincoKgFourZeroRThreeZeroMainFpr\%にとどまり，$k_g{=}3$は$r{=}1$で\NKgNincoKgThreeROneMainFpr\%と最も悪いが$r{=}30$で\NKgNincoKgThreeRThreeZeroMainFpr\%まで下がる．$\sigma$が静的$p$を初めて下回る$r$は$k_g\in\{3,5,10,20\}$で\NKgNincoKgOneZeroCrossR{}，$k_g{=}40$で\NKgNincoKgFourZeroCrossR{}である．SSB-hard（b）では$k_g\ge20$は$r{=}1$から，$k_g\le10$は$r{=}2$から静的$p$を下回り，$r{=}50$では$k_g{=}5$（\NKgSsbHardKgFiveRmaxMainFpr\%）と$10$（\NKgSsbHardKgOneZeroRmaxMainFpr\%）が最も良く，$k_g{=}40$は\NKgSsbHardKgFourZeroRmaxMainFpr\%である．$\pP$だけ（破線）も同じ順に並び，種の利得（実線と破線の差）は$r\ge10$で現れ，$k_g{=}40$では小さい．密なグラフでは再出現する未知の塊の内側の辺の割合が下がり，塊が孤立しにくい（命題2のカット条件）と解釈できる．登録時の仮説H5（付録\ref{sec:app_hyp}）の前半（小さい$k_g$ほど小さい$r$で静的$p$を下回る）は否定され，後半（大きい$r$での最良は$k_g\in\{10,20\}$）は$\{5,10\}$として部分的に成り立つ．本稿の$k_g{=}10$は登録どおり変えない．
